% Sección: Resultados
%
% ¿Valdrá la pena mostrar algunos resultados intermedios, anteriores a la
% evaluación en los datos de testeo? Umbrales de decisión, métricas vs.
% hiperparámetros, etc.

\subsection{Desempeño de los modelos implementados}

% Describir verbalmente el desempeño y mostrar métricas calculadas sobre los
% datos de testeo. ¿Mostrar gráficos, modelos y métricas por modelo? ¿Mostrar
% cuadros de métricas por modelo? Hacer una breve interpretación de los
% resultados en relación con la temática.

% Métricas de los modelos calculadas en test
\begin{table}[h]
\begin{tabular*}{\textwidth}{@{\extracolsep{\fill}}lrrrrr}
\toprule
  Modelo & $F_2$ & \textit{Accuracy} & \textit{Precision} & \textit{Recall} & \textit{AUC (PR)}\\
\midrule
Random Forest & 0.9941 & 0.9883 & 0.9710 & 1.0000 & 0.9993\\
Logit LASSO & 0.9911 & 0.9825 & 0.9571 & 1.0000 & 0.9993\\
KNN & 0.9853 & 0.9708 & 0.9306 & 1.0000 & 1.0000\\
Logit LASSO (PCA) & 0.9851 & 0.9883 & 0.9851 & 0.9851 & 0.9973\\
LDA & 0.9735 & 0.9649 & 0.9296 & 0.9851 & 0.9961\\
Logit & 0.9552 & 0.9649 & 0.9552 & 0.9552 & 0.9891\\
\bottomrule
\end{tabular*}

\caption{Métricas de los modelos en los datos de testeo. Se presentan en orden descendiente según el puntaje $F_2$ obtenido en el conjunto de testeo.}
\label{tab:metricas}
\end{table}

En el Cuadro \ref{tab:metricas} se presentan las métricas obtenidas por cada
modelo sobre los datos de testeo. Los dos modelos que alcanzaron los mayores
valores en $F_2$ fueron el de bosques aleatorios con 0.9941, y la regresión
logística con penalización LASSO sobre las variables originales con 0.9911. En
ambos casos la clasificación fue casi perfecta: los bosques aleatorios lograron
una precisión de 0.971 con sólo dos falsos positivos y un \textit{recall}
perfecto; la regresión logró una precisión de 0.957 con tres falsos positivos
y, también, \textit{recall} perfecto. El resto de los modelos no se quedó
demasiado atrás: el que tuvo el peor desempeño fue la regresión logística sobre
las componentes principales de las variables originales con un puntaje
$F_2$ de 0.955 y seis tumores mal clasificados, tres benignos y tres malignos.

% Matrices de confusión de los modelos elegidos
\begin{table}[h]
\begin{tabular*}{\textwidth}{@{\extracolsep{\fill}}lcccc}
\toprule
\multirow{3}{3em}{Realidad}  & \multicolumn{4}{c}{Predicción} \\
          & \multicolumn{2}{c}{Bosques aleatorios} & \multicolumn{2}{c}{Logística LASSO} \\
            \cmidrule{2-5}
          & Maligno & Benigno  & Maligno & Benigno \\
\midrule
Maligno   & 67      & 0        & 67      & 0 \\
Benigno   & 2       & 102      & 3       & 101 \\
\bottomrule
\end{tabular*}

\caption{Matrices de confusión de los dos mejores modelos.}
\label{tab:matrices}
\end{table}

Es interesante notar el efecto de la ponderación en el puntaje $F_\beta$ sobre
el resultado. La regresión logística con penalización LASSO sobre las
componentes principales de las predictoras tuvo la misma \textit{accuracy} que
el modelo de bosques aleatorios (dos tumores mal clasificados). Sin embargo,
como se puede ver en el Cuadro \ref{tab:matrices}, una de sus predicciones fue
un falso positivo, por lo cual terminó teniendo un puntaje $F_2$ menor al de
los bosques aleatorios.

\subsection{Descripción de los modelos seleccionados}\label{sec:detalle}

% Describir más detalladamente el procedimiento de cada modelo. ¿Quizás aquí van
% los gráficos de resultados intermedios? Mostrar los parámetros, los del pca y
% los de las regresiones si es que fueran seleccionadas. ¿Mencionar lo de
% tidymodels acá?

\subsubsection{Bosques aleatorios}

Para entrenar este modelo se utilizaron las variables originales normalizadas
salvo las asociadas a la dimensión fractal (\verb|fractal_dimension\_mean|,
\verb|..._se| y \verb|..._worst|) y el desvío estándar de la textura
(\verb|texture_se|), la suavidad (\verb|smoothness_se|) y la simetría
(\verb|symmetry_se|), que se excluyeron dado que en los histogramas de densidad
no se vió que ofrecieran capacidad de separar las dos clases de tumores. Como
los árboles aleatorios utilizan una única variable por \textit{split}, no tiene
ventaja incluir aquellas cuya distribución es similar en ambas clases. Por un
motivo similar se decidió no buscar las componentes principales de las
variables; la alta correlación entre las predictoras no debería entrar en juego
si en cada \textit{split} sólo se usa una de ellas.

Para el modelo en sí se utilizó la implementación del paquete \verb|{ranger}| a
través de \verb|{tidymodels}|. Los hiperparámetros de este modelo son los
siguientes: el número de árboles a entrenar, que, tras notar que no producía
cambios notorios en la métrica de validación, se estableció en 100, un número
relativamente bajo para no aumentar innecesariamente el costo computacional; el
número de observaciones mínimo por hoja, que se buscó optimizar; y el número de
variables a considerar en cada split de cada árbol, que también se optimizó. En
este proceso se consideraron los valores de 10, 20 y 30 observaciones mínimas
por hoja, y entre 1 y 10 variables por split, y se buscó sobre una grilla con
todas las combinaciones de estos dos grupos de valores. Como se mencionó antes,
la optimización se hizo mediante validación cruzada de cinco \textit{folds}
maximizando el área bajo la curva \textit{precision-recall}. Los \textit{folds}
se crearon estratificando por \verb|target| para mantener la proporción de
casos malignos dentro de cada uno.

% Métricas de validación de random forest
\begin{figure}[h]
\includegraphics[width=\textwidth]{figuras/rf\_metricas.pdf}
\caption{Métricas de desempeño del modelo de bosques aleatorios según hiperparámetros optimizados. Los gráficos están segmentados según el número mínimo de observaciones por hoja.}
\label{fig:rfmetricas}
\end{figure}

En la Figura \ref{fig:rfmetricas} se puede ver la métrica de validación
calculada según los hiperparámetros junto con la precisión y el \textit{recall}
calculados con un umbral de 0.5. La combinación que rindió la mayor área bajo
la curva fue una variable por split y diez observaciones mínimo por hoja de
cada árbol con un área bajo la curva promedio de 0.985 y error estándar de
0.004 aproximadamente. Algo que se observa es que el área no varía demasiado
según el hiperparámetro, pero sí lo hacen la precisión y el \textit{recall}.

% Acumulada de las probabilidades predichas de random forest
\begin{figure}[h]
  \includegraphics[width=\textwidth]{figuras/rf\_pred\_umbral.pdf}
  \caption{(A) Distribución acumulada de las probabilidades predichas y (B) puntaje $F_2$ según umbral de decisión correspondientes al modelo de bosques aleatorios. La línea punteada vertical indica el umbral de decisión elegido.}
\label{fig:rfpredumb}
\end{figure}

Una vez obtenidos los hiperparámetros optimizados se utilizaron las
predicciones hechas sobre los cinco \textit{folds} para buscar el umbral de
decisión óptimo que maximizase el puntaje $F_2$. Esto se hizo sobre el total de
las predicciones sin distinguir el \textit{fold} al que perteneciesen. Se buscó
sobre una grilla de 99 valores equiespaciados entre 0.01 y 0.99, y se tomó como
óptimo el promedio de los umbrales que alcanzasen el máximo $F_2$ (ya que podía
haber más de un valor que tocase el máximo). En la Figura \ref{fig:rfpredumb}(A) se
muestran las distribuciones acumuladas de las probabilidades predichas de ser
un tumor maligno según la condición real del tumor y el \textit{fold} de
validación. Se ve que el modelo logra una buena separación de los casos, pero
pareciera que en una buena proporción de los casos malignos las probabilidades
predichas son demasiado bajas; en los cinco \textit{folds} aproximadamente el
12.5\% tiene una probabilidad predicha menor a 0.5. La figura
\ref{fig:rfpredumb}(B) muestra el puntaje $F_2$ alcanzado según el umbral de
decisión y \textit{fold}. Se puede ver que el máximo suele alcanzarse en
valores más bajos que 0.5. La línea negra punteada se ubica en 0.29 en el eje
$x$ e indica el mejor umbral encontrado, donde se alcanza un $F_2$ de 0.9435.

En el Cuadro \ref{fig:rfespec} del Anexo \ref{sec:anexo1} se resumen las características del modelo final de forma más sintética.


\subsubsection{Regresión logística con penalización LASSO}

Para este modelo no se realizó una selección de variables ya que la
penalización LASSO tiene como consecuencia que algunos de los coeficientes
estimados de la regresión pueden ser exactamente 0. Se utilizaron todas las
variables predictoras disponibles en el conjunto de datos y se las normalizó
para poder implementar el modelo. El término de penalización LASSO es $\lambda
\sum_{j=1}^p | \beta_j |$, donde $p$ es la cantidad de variables predictoras,
en este caso 30, y $\beta_j$ es el coeficiente correspondiente a la predictora
$j$. Como los coeficientes $\beta_j$ están en la unidad de sus respectivas
variables, si estas no se normalizaran aquellas con unidades más grandes (como
las relacionadas con \verb|area| en nuestro caso) tendrían más peso en el
término y serían las primeras en perder influencia. Si bien el paquete
\{glmnet\} normaliza las variables de forma automática, se optó por dejar por
explícito este paso del preprocesamiento de los datos en el código de R.

Se utilizó la implementación del paquete \verb|{glmnet}| mediante
\verb|{tidymodels}| de R. Este paquete, en realidad, ofrece modelos lineales
generalizados con penalización \textit{elastic net}, que es una combinación
lineal de la penalización LASSO y \textit{Ridge}
\parencite{tay_etal_2023_elastic}. El hiperparámetro que controla esa
combinación es $\alpha$ (\verb|mixture| en la implementación usada), que se
fijó en 1 para utilizar exclusivamente la penalización LASSO. Esto se hizo, por
un lado, para simplificar la búsqueda de hiperparámetros óptimos y, por otro
lado, porque es razonable que en este escenario sea útil seleccionar variables
predictoras para luego priorizar qué aspectos de las células de los tumores
medir. La búsqueda se hizo de forma análoga al modelo de bosques aleatorios; se
creó una grilla de 25 valores equiespaciados en escala logarítmica de base 10
entre $10^{-6}$ y $10^{-1}$ y se seleccionó el que resultase en la mayor área
bajo la curva \textit{precision-recall}. Se utilizaron los mismos
\textit{folds} que en la validación del modelo de bosques aleatorios.

% Métricas de validación de lasso
\begin{figure}[h]
\includegraphics[width=\textwidth]{figuras/lasso1\_metricas.pdf}
\caption{Métricas de desempeño de la regresión logística con penalización LASSO según parámetro de penalización $\lambda$.}
\label{fig:lassometricas}
\end{figure}

La Figura \ref{fig:lassometricas} muestra el valor de las mismas tres métricas
que en el modelo anterior según el valor del hiperparámetro $\lambda$. El mejor
valor encontrado fue $\lambda$ = 0.0056 con un área de 0.991 y error estándar
de 0.0031 aproximadamente. Se ve que cuando crece el coeficiente de
penalización la precisión (con umbral 0.5) sube y el \textit{recall} cae,
seguramente porque el modelo está prediciendo mayormente negativos, lo cual
reduce los falsos positivos y aumenta los falsos negativos. Para que se dé este
resultado, de todas formas, los verdaderos positivos no deberían bajar al mismo
ritmo que los falsos positivos, caso contrario la precisión no debería cambiar
mucho.

% Acumulada de las probabilidades predichas de lasso y umbral vs F2 
\begin{figure}[h]
  \includegraphics[width=\textwidth]{figuras/lasso1\_pred\_umbral.pdf}
  \caption{(A) Distribución acumulada de las probabilidades predichas y (B) puntaje $F_2$ según umbral de decisión correspondientes al modelo de regresión logística con penalización LASSO. La línea punteada vertical indica el umbral de decisión elegido.}
\label{fig:lassopredumb}
\end{figure}

La búsqueda del umbral óptimo se realizó de la misma manera que en el modelo
anterior. La Figura \ref{fig:lassopredumb} es análoga a la figura
\ref{fig:rfpredumb}. Se ve que este modelo también logra una buena separación
de los casos y, en comparación con el anterior, asigna probabilidades más altas
a los tumores malignos, con alrededor del 70\% de las observaciones malignas
con probabilidad cercana a 1. Sin embargo, sigue habiendo casos de este tipo
con predicciones más bajas de lo deseable que se asemejan a las de los casos
benignos con predicciones más altas. El umbral que maximiza el puntaje $F_2$
(calculado sobre el total de predicciones sobre los \textit{folds}) es 0.2 con
un valor $F_2$ de 0.9595.

El modelo final estimó 12 parámetros distintos de cero, incluyendo el
intercepto, que se listan en el Cuadro \ref{fig:lassoespec}, donde se resumen las
especificaciones. Sorprendentemente, la mayor parte de las variables
seleccionadas fueron los peores casos (\verb|_worst|) de ciertas variables,
específicamente, del radio, la textura, el perímetro, la suavidad, la
concavidad, la cantidad de secciones cóncavas y la simetría. Fueron
seleccionadas tres variables de desvío estándar (\verb|_se|): la del radio, la
compacidad y la dimensión fractal. Sólo una corresponde a la media
(\verb|_mean|), la del número de secciones cóncavas. La mayoría de las
variables aumentan las chances de que el tumor sea maligno; sólo el desvío de
la compacidad y de la dimensión fractal reducen esas chances.
